    \subsection{PRG for low-communication functions in the number-in-hand setting}\label{appendix:prg_cc}
	
	In this subsection, we show how to fool functions with low communication complexity in the number-in-hand model.
    
	\begin{theorem}[\cite{ASWZ96,INW94}, \Cref{thm:prg-lc} restated]\label{lem:prg-lc-app}
		For any $k\geq 2$, there exists a PRG that $\delta$-fools any $n$-bits functions with $k$-party number-in-hand deterministic communication complexity at most $D'$, with seed length
		\[
			n/k+O\left(D' + \log(1/\delta)+\log(k)\right)\cdot\log(k).
		\]
	\end{theorem}
	
	The PRG in \Cref{thm:prg-lc} is based on the PRG by Impagliazzo, Nisan and Wigderson~\cite{INW94} that is used to derandomize ``network algorithms" and space-bounded computation. We will need to use randomness extractors, which we review below.
	
	\begin{definition}[Min-entropy]
    	Let $X$ be a random variable. The \emph{min-entropy} of $X$, denoted by $H_\infty(X)$, is the largest real number $k$ such that $\Prob\csqbra{X=x}\leq2^{-k}$ for every $x$ in the range of $X$. If $X$ is a distribution over $\boolpm^\EI$ with $H_\infty\cprn{X}\geq k$, then $X$ is called a {\em$\prn{\EI,k}$-source}.
    \end{definition}
    
    \begin{definition}[Extractors]
    	A function $\mathrm{Ext}\colon\boolpm^\EI\times\boolpm^d\to\boolpm^m$ is an \emph{$(k,\varepsilon)$-extractor} if, for any $\prn{\EI,k}$-source $X$, and any test $T\colon\boolpm^m\to\boolpm$, it is the case that
    	\[
    	    \left|\Prob[T(\mathrm{Ext}\cprn{X,\uniform_d}X)=1]-\Prob[T(\uniform_m)=1]\right|\leq \varepsilon.
    	\]
    \end{definition}
	
    \begin{theorem}[{\cite[Theorem 6.22]{Vad12}}]\label{thm:prg-ext}
		For any integer $m,\kappa>0$ and $0<\delta'<0$, there exists an explicit $(\kappa,\delta')$ extractor $\mathrm{Ext}\colon\bool^{m}\times\bool^d\to\bool^{m}$ with $d=O(m-k+\log(1/\delta'))$.
	\end{theorem}

    We are now ready to show \Cref{thm:prg-lc}.
    
	\begin{proof}[Proof of \Cref{thm:prg-lc}]\renewcommand{\qedsymbol}{$\square$ (\Cref{thm:prg-lc})}
	    We first describe the construction of the PRG. In fact, we will construct a sequence of PRGs $G_0,G_1,\dots,G_{\log(k)}$. We begin by specifying the parameters of these PRGs.
	    Let $t=\log(k)$, and let
	    \[
	    	d=O\left(D'+\log(1/\delta)+t\right).
	    \]
	    For $i=0,1,\dots,t$, let
	    \begin{itemize}
	        \item $r_0=n/k$,
	        \item $r_i=r_{i-1}+d$.
	    \end{itemize}
	    Note that we have $r_i=n/k+i\cdot d$. Also, let
	    \[
	    	\mathrm{Ext}_i\colon\bool^{r_i}\times\bool^d\to\bool^{r_i}
	    \]
	    be a $(\kappa_i,\delta')$-extractor from \Cref{thm:prg-ext}, where
	    \[
	    	\kappa_i=r_i-D'-2t - \log(1/\delta)
	    \] 
	    and
	    \[
	        \delta'=\delta/\left(3^t\cdot 2^{D'}\right).
	    \]
	    Note that the seed length of the extractors is $d=O\left(D'+\log(1/\delta)+t\right)$. Finally, define $G_i\colon\bool^{r_i}\to\bool^{n/2^{t-i}}$ recursively as follows
	    \begin{itemize}
	    	\item $G_{0}(a)=a$, where $a\in\bool^{n/k}$.
	    	\item $G_i(a,z)=G_{i-1}(a)\circ G_{i-1}(\mathrm{Ext}_{i-1}(a,z))$, where $a\in\bool^{r_{i-1}}$ and $z\in\bool^d$.
	    \end{itemize}
	    We will show that $G_t\colon\bool^{r_t=n/k+t\cdot d}\to\bool^n$ fools any functions $f$ with $k$-party number-in-hand deterministic communication complexity at most $D'$.
	    First, note that such $f$ can be written as
	    \[
	    	f(x_1,x_2,\dots,x_k)=\sum_{i=1}^{2^{D'}} h^{(i)}_{1}(x_1)\cdot h^{(i)}_{2}(x_2)\cdot\ldots\cdot h^{(i)}_{k}(x_k),
	    \]
	    for some $h^{(i)}_{j}\colon\bool^{n/k}\to \bool$ ($i\in\left[2^{D'}\right], j\in[k]$).
	    Therefore, to show that the PRG $G_t$ $\delta$-fool $f$, it suffices to show that $G_t$ $\left(\delta/2^{D'}\right)$-fools every function $g$ of the form
	    \[
	    	g(x_1,x_2,\dots,x_k)=h_1(x_1)\cdot h_2(x_2)\cdot\ldots\cdot h_k(x_k).
	    \]
	    More specifically we show the following.
	    
	    \begin{claim}\label{lem:prg-lc-claim-1}
	    	For every $k \geq 2$ and $  0\leq i \leq t$, the generator $G_i$ defined above $\left(3^{i}\cdot \delta'\right)$-fools every function $g_i\colon\bool^{n/2^{t-i}}\to\bool$ of the form
	    	\[
	    		g_i(x_1,x_2,\dots,x_{k/2^{t-i}})=h_1(x_1)\cdot h_2(x_2)\cdot\ldots\cdot h_{k/2^{t-i}}(x_{k/2^{t-i}}),
	    	\]
	    	where $x_1,x_2,\dots,x_{k/2^{t-i}}\in\bool^{n/k}$.
	    \end{claim}
	    
	    \begin{proof}\renewcommand{\qedsymbol}{$\square$ (\Cref{lem:prg-lc-claim-1})}
	    	The proof is by induction on $i$. The base case is $i=0$, which is trivial given the definition of $G_{0}$. Now suppose the claim holds for $i-1$, we show the case for $i$. This is done using a hybrid argument. Consider the following four distributions
	    	\begin{itemize}
	    		\item $\mathcal{D}_1=U_{n/2^{t-i}}$,
	    		\item $\mathcal{D}_2=U_{n/2^{t-i+1}}\circ G_{i-1}(U_{r_{i-1}})$,
	    		\item $\mathcal{D}_3=G_{i-1}(U_{r_{i-1}}) \circ G_{i-1}(U'_{r_{i-1}})$ ($U$ and $U'$ are two independent uniform distributions),
	    		\item $\mathcal{D}_4=G_i(U_{r_i})$.
	    	\end{itemize}
	    	We want show show that
	    	\[
	    		\left|\Exp[g_i(\mathcal{D}_1)]-\Exp[g_i(\mathcal{D}_4)]\right|\leq 3^{i}\cdot \delta'.
	    	\]
	    	By the triangle inequality, it suffices to show that
	    	\begin{equation}\label{eq:prg-lc-triangle}
	    		\left|\Exp[g_i(\mathcal{D}_1)]-\Exp[g_i(\mathcal{D}_2)]\right| + \left|\Exp[g_i(\mathcal{D}_2)]-\Exp[g_i(\mathcal{D}_3)]\right| + \left|\Exp[g_i(\mathcal{D}_3)]-\Exp[g_i(\mathcal{D}_4)]\right| \leq  3^{i}\cdot \delta'.
	    	\end{equation}
	    	We show \Cref{eq:prg-lc-triangle} by upper bounding each of the three summands.
	    	
	    	\paragraph{First summand.}
	    	
	    	We show that
            \begin{equation}\label{eq:prg-lc-triangle-1}
				\left|\Exp[g_i(\mathcal{D}_1)]-\Exp[g_i(\mathcal{D}_2)]\right| \leq 3^{i-1}\cdot \delta'.
			\end{equation}
			Let us re-write $g_i$ as
	    	\[
	    		g_i(x_1,x_2,\dots,x_{k/2^{t-i}})=\LH{h}(x_1,x_2,\dots,x_{k/2^{t-i+1}})\cdot \RH{h}(x_{k/2^{t-i+1}+1},x_{k/2^{t-i+1}+2},\dots,x_{k/2^{t-i}}),
	    	\]
	    	where
	    	\[
	    	    \LH{h}(y)\vcentcolon=\prod_{j=1}^{k/2^{t-i+1}}h_i(y) \text{\quad and\quad  } \RH{h}(y)\vcentcolon=\prod_{j=k/2^{t-i+1}}^{k/2^{t-1}}h_i(y).
	    	\]
	    	Then,
	    	\begin{align*}
	    		\Exp[g_i(\mathcal{D}_2)]&=\Exp\left[\LH{h}(U_{n/2^{t-i+1}})\cdot \RH{h}(G_{i-1}(U_{r_{i-1}}))\right]\\
	    		&=\Exp\left[\LH{h}(U_{n/2^{t-i+1}})\right]\cdot\Exp\left[\RH{h}(G_{i-1}(U_{r_{i-1}}))\right]\\
	    		&=\Exp\left[\LH{h}(U_{n/2^{t-i+1}})\right]\cdot\left(\Exp\left[\RH{h}(U_{n/2^{t-i+1}})\right]\pm 3^{i-1}\cdot \delta'\right)
	    		\tag{By the induction hypothesis}\\
	    		&=\Exp\left[\LH{h}(U_{n/2^{t-i+1}})\right]\cdot\Exp\left[\RH{h}(U_{n/2^{t-i+1}})\right] \pm 3^{i-1}\cdot \delta'\\
	    		&=\Exp[g_i(\mathcal{D}_1)]\pm 3^{i-1}\cdot \delta',
	    	\end{align*}
	    	as desired.
	    	
	    	\paragraph{Second summand.}
	    	
	    	By a similar argument, it can be shown that
            \begin{equation}\label{eq:prg-lc-triangle-2}
				\left|\Exp[g_i(\mathcal{D}_2)]-\Exp[g_i(\mathcal{D}_3)]\right| \leq 3^{i-1}\cdot \delta'.
			\end{equation}
			We omit the details here.
	    	
	    	\paragraph{Third summand.}
	    	
	    	We show that
            \begin{equation}\label{eq:prg-lc-triangle-3}
				\left|\Exp[g_i(\mathcal{D}_3)]-\Exp[g_i(\mathcal{D}_4)]\right| \leq \delta'.
			\end{equation}
	    	We have
	    	\begin{align*}
	    		\Exp[g_i(\mathcal{D}_4)]&=\Exp[g_i(G_i(U_{r_i}))]\\
	    		&=\Exp\left[\LH{h}(G_{i-1}(X))\cdot \RH{h}(G_{i-1}(\mathrm{Ext}_{i-1}(X,Z)))\right]
	    		\tag{where $X\sim\bool^{r_{i-1}}$ and $Z\sim\bool^{d}$}\\
	    		&=\Exp[A(X)\cdot B(\mathrm{Ext}_{i-1}(X,Z))]
	    		\tag{where $A(\cdot)= \LH{h}(G_{i-1}(\cdot))$  and $B(\cdot)= \RH{h}(G_{i-1}(\cdot))$}\\
	    		&=\Exp[B(\mathrm{Ext}_{i-1}(X,Z))\given A(X)=1]\cdot\Prob[A(X)=1].
	    	\end{align*}
	    	Similarly, we get
	    	\[
	    		\Exp[g_i(\mathcal{D}_3)]=\Exp[B(U_{r_{i-1}})\given A(X)=1]\cdot\Prob[A(X)=1].
	    	\]
	    	As a result, we have
	    	\begin{align*}\label{eq:prg-lc-triangle-3.1}
	    		&\quad\left|\Exp[g_i(\mathcal{D}_4)]-\Exp[g_i(\mathcal{D}_3)]\right|\\
	    		&=\left|\left(\Exp[B(\mathrm{Ext}_{i-1}(X,Z))\given A(X)=1]-\Exp[B(U_{r_{i-1}})\given A(X)=1]\right)\cdot\Prob[A(X)=1]\right|.
	    		\numberthis
	    	\end{align*}
		    On the one hand, if $\Prob[A(X)=1]\leq \delta'$, then \Cref{eq:prg-lc-triangle-3.1} is at most $\delta'$. On the other hand, if $\Prob[A(X)=1]> \delta'$, then
	        \[
	        	H_{\infty}\cprn{X\given A(X)=1} > r_{i-1}-\log(1/\delta')> r_{i-1}-D'-2t -\log(1/\delta)=\kappa_{i-1}.
	        \]
	        Then by the fact that $\mathrm{Ext}_{i-1}$ is a $(\kappa_{i-1},\delta')$-extractor, we have
	        \[
	            \left|\Exp[B(\mathrm{Ext}_{i-1}(X,Z))\given A(X)=1] - \Exp[B(U_{r_{i-1}})\given A(X)=1]\right| \leq \delta'.
	        \]
	        Therefore, \Cref{eq:prg-lc-triangle-3.1} is at most $\delta'$ and this complete the proof of \Cref{eq:prg-lc-triangle-3}. Finally, note that \Cref{eq:prg-lc-triangle} follows from \Cref{eq:prg-lc-triangle-1}, \Cref{eq:prg-lc-triangle-2} and \Cref{eq:prg-lc-triangle-3}. This completes the proof of \Cref{lem:prg-lc-claim-1}.
	    \end{proof}
	    
	    Given \Cref{lem:prg-lc-claim-1}, \Cref{thm:prg-lc} now follows by letting $i=t$.
	\end{proof}
